\subsection{L 集群环境细节}\label{l-cluster-environment-details}

本附录介绍使训练环境在前沿（frontier）规模下可用的基础设施机制。核心设计选择是将物理拓扑与硬件健康状态作为一等（first-class）调度状态暴露给调度系统。节点与 NVLink 机架域通过拓扑标签（topology labels）、预留（reservation）对象、机架级服务和认证门控（certification gates）来表示。一个 node 并非仅仅因为已被供应（provisioned）就视为可用；只有当它健康、拓扑有效、可观测且可恢复时，才算真正可用。

\subsection{L.1 硬件与计算集群}\label{l.1-hardware-and-compute-clusters}

MAI-Thinking-1 是在一个异构加速器机群上开发的，涵盖 NVIDIA H100、GB200 和 GB300 系统。不过，主预训练运行被放置在单个站点的单一 GB200 集群上。GB200 和 GB300 集群托管在 Microsoft 自营数据中心内，并通过与 Azure 团队联合开发的自定义镜像提供给 MAI。这些系统按机架级 NVL72 单元供应，每个机架提供 72-GPU NVLink 域，用于高带宽的 scale-up（纵向扩展）通信；机架间 RDMA 使用 InfiniBand。H100 系统也仍保留在实验室环境中，使用 8-GPU node，配备节点内 NVLink/NVSwitch 和 InfiniBand 用于 scale-out（横向扩展）通信。

逻辑集群供应与管理。每个站点被划分为多个 Kubernetes 集群，通常每个数据中心楼栋一个。逻辑集群包含用于训练的 GPU node 和用于支撑服务的 CPU node。自定义控制器（custom controllers）将期望的逻辑集群状态与底层 Azure 资源调和（reconcile），维护节点拓扑标签，并将外部健康数据库集成到调度状态中。

拓扑标签在 node 启动时根据物理主机身份分配。这些标签表示机架、主机和位置（locality）信息，供调度、预留和机架级服务使用。这使得物理放置对控制平面（control plane）可见，让作业能够落在不仅可用、而且拓扑上合适的容量上。

大型作业在需要时被设计为跨越单个 Kubernetes 集群的边界。node 在整个计算环境中均可全局路由，工作负载 pod 使用 host networking，以避免不必要的 overlay 开销。

\subsection{L.2 集群就绪与认证}\label{l.2-cluster-readiness-and-certification}

在数千 GPU 规模下，故障是预期之中的。因此认证（certification）充当第一道可靠性边界。其目的是防止坏节点、退化链路、边缘存储和静默损坏（silent-corruption）风险进入生产训练池。认证框架是自动化的、Kubernetes 原生（Kubernetes-native）的且分层的。硬件必须先通过逐步更广范围的测试，然后才可调度：单节点诊断、机架级多节点集合通信（collectives）以及选定的跨机架 InfiniBand 验证。新节点和修复后的节点遵循相同的路径。从维修返回的 node 必须再次通过认证才能重新进入生产，防止重复出错的节点（repeat offenders）循环回到大型训练作业中。

认证。我们的诊断套件对每个节点进行独立验证。与通用健康检查不同，这些测试会对单个组件进行长时间压力测试：GPU、CPU 核、HCA、NVLink 链路和主内存，从而在 node 被接纳进生产之前，隔离出高 ECC 率、异常节流行为、GPU 时钟违规、链路抖动（link flaps）及其他边缘状况。

该套件还使用 NCCL 集合通信对每个机架的节点内和节点间通信路径进行压力测试，以识别异常的 NVLink 行为。这些自动化测试在稳定 GB200 和 GB300 机架方面发挥了关键作用：我们识别出多个健康节点少于 16 个的机架，这些机架被排除在大型训练运行之外。

认证的最后阶段包括跨机架 InfiniBand 测试，覆盖单节点和机架级测试无法触及的 fabric 路径。这些测试将跨机架、rail 和 leaf 组的节点分组，以验证 spine 层行为、路径多样性和跨机架 RDMA 性能。由于它们的 blast radius（影响范围）更大，因此在 bring-up（启动）、事故调查和 fabric 验证期间按需运行。

节点生命周期。认证是闭环节点生命周期中的一个阶段，如图 26 所示。新的或修复后的 node 进入 Init，在那里被打上 taint、完成初始化、标记 label 并配置。随后，认证控制器运行分层测试套件。通过的节点变为 Available，失败的节点变为 Impaired。当检测到 NPD 条件、XID 错误、ECC 阈值、NVLink 退化、InfiniBand 链路抖动或存储故障时，运行时监控也可将 Available 节点移至 Impaired。瞬态问题（如重启或 soft drain（软排空））进入 Auto Remediation（自动修复）。持续性或归因于硬件的问题通过 Guest Health Reporting 路由至厂商或数据中心维护。修复后的节点返回 Init 并重新进入认证流程。无法恢复的节点被退役（decommissioned）。

\subsection{L.3 调度、编排与控制平面}\label{l.3-scheduling-orchestration-and-control-plane}

调度问题是在异构容量上放置异构工作负载，同时不牺牲局部性（locality）、配额隔离或恢复能力。工作负载包括长时间预训练作业、后训练和强化学习运行、推理与评测作业，以及 CPU 密集型数据管线。因此技术栈是分层的：Kubernetes 维护集群状态，Kueue 做出 admission 和配额决策，MAI 控制器维护预留与拓扑就绪，Ray 执行分布式作业，而机群范围（fleet-wide）的控制平面统一了跨集群和调度后端的可见性。

\pandocbounded{\includegraphics[keepaspectratio,alt={image}]{images/f2cd7e527ffef59be6a85867920a4cd64e7eda2c4b2527677bc9f80cbeebad98.jpg}}

图 26. 节点生命周期：从初始化到认证、生产使用、修复、维护与退役。

集群本地控制平面。集群本地的 MAI 控制平面管理预留、机架拓扑、配额一致性和调度就绪门控。它不替代调度器，而是维护调度器所需的状态。Kueue 处理优先级、配额、admission、preemption（抢占）和拓扑感知放置。当一个工作负载被接纳（admitted）时，Kueue 会尽量保持高带宽域内的局部性，使大型作业落在紧凑的拓扑区域，而不是零散的容量上。机架碎片化（rack fragmentation）是反复出现的调度风险。来自许多队列的小作业可能会填满机架，从而阻止后来的大型作业获取连续容量。集群本地控制平面通过软机架预留（soft rack reservations）解决这一问题。队列被分配优先机架，工作负载会优先落在这些机架上。当容量空闲时仍可借用（borrowing），但 reclaimWithinCohort 允许队列在需要时收回其预留机架。这提高了负载下的装箱（bin packing）效率，同时在低需求时保持利用率。

Ray 运行时。Ray 在被接纳的作业内部提供分布式运行时。MAI 驱动程序将已接纳的拓扑转换为 actor 放置、通信组和 NCCL clique 配置。预训练作业要求严格的 learner 可用性，而 RL 作业管理多种 actor 类型，包括 learner、推理服务器、rollout worker 和路由器，每种都有不同的放置和容错需求。驱动程序监控 actor 存活状态、协调训练循环，并在异步组件之间保持 checkpoint 一致性。

\subsection{L.4 可观测性、遥测与集群监控}\label{l.4-observability-telemetry-and-fleet-monitoring}

我们将底层硬件遥测、调度器状态、作业元数据和机群效率（fleet-efficiency）信号整合到统一的运维视图中。GPU 健康状态通过 XID、ECC、热状态、功耗、时钟节流、NVLink 状态、NVLink 误码率、chip-to-chip 链路、InfiniBand 设备状态、本地 NVMe 健康、PCIe 错误和驱动状态进行跟踪。这些信号被转换为 Kubernetes node conditions（节点条件），再通过自定义 triage 和 drain 控制器转化为调度与修复动作。这使得健康状态不仅对运维人员可观测，对调度器也可观测。

作业可观测性聚合自 Kueue、Kubernetes、Ray、训练日志和实验元数据。运维人员可以按 namespace、pod、job 和重启索引（restart index）检查队列、优先级、admission 状态、节点放置、worker 就绪状态、重启次数、训练配置、step 进度以及作用域日志（scoped logs）。这使得无需手动跨系统重建状态，即可区分调度延迟、运行时故障、节点故障、存储退化和应用级停滞（application-level stalls）。

遥测数据存储并查询于针对不同时间尺度（time horizons）优化的系统上。Datadog 提供近实时指标和日志搜索。Azure Managed Prometheus 提供集群内和跨集群（cross-cluster）的时间序列采集。Azure Data Explorer 提供长期保留的分析存储，用于日志、指标、存储遥测和集群状态。告警同样采用分层模型：本地 Prometheus 规则用于低延迟组件故障，Datadog monitor 用于服务和日志衍生信号，基于 KQL 的检查用于长期保留分析，Azure Monitor 用于资源和 Prometheus 告警。关键告警被路由到事件管理（incident management），而较低严重程度的信号则流入仪表盘和运维评审。

独特的设计之处在于，可观测性是控制回路的一部分。硬件遥测、fabric 健康、存储行为、调度状态和作业进度并不是被动的仪表盘；它们决定容量是被接纳、排空、修复，还是重新投入服务。这使得整个机群可以按可用训练容量和有效吞吐（goodput）来管理，而不仅仅是按已供应的 GPU 数量。
